% =====================================================================
%  CHAPTER 11: 学校突发公共卫生事件应急管理（对应大纲第十一章）
%  内容来源：往届笔记第十六章，参考PPT第十八章
% =====================================================================
\chapter{学校突发公共卫生事件应急管理}
\setcounter{chapter}{11}
\chapterintro{
\textbf{掌握：}学校突发公共卫生事件概述。学校传染病事件、食物中毒、群体心因性反应事件的应对。
}

% ===== 掌握内容 =====
\section{学校突发公共卫生事件概述}

\subsection{定义与类型}

\begin{defbox}
\textbf{学校突发公共卫生事件（public health emergency events in schools）：}在学校内突然发生，造成或可能造成师生员工健康严重损害的重大传染病疫情、群体性不明原因疾病、重大食物或职业中毒及其他严重影响师生员工身心健康的公共卫生事件。

\textbf{类型（7类）：}

\begin{enumerate}[leftmargin=*]
    \item \imp{重大传染病疫情}
    \item \imp{预防接种和预防服药群体性不良反应}
    \item \imp{群体性不明原因疾病}
    \item \imp{食物中毒}
    \item \imp{其他中毒}
    \item \imp{环境因素事件}
    \item \imp{意外辐射照射事件}
\end{enumerate}
\end{defbox}

\subsection{分级}

\begin{keybox}
\textbf{学校突发公共卫生事件分级（从高到低）：}

\begin{table}[H]
\centering
\caption{学校突发公共卫生事件分级}
\begin{tabularx}{\textwidth}{l>{\centering\arraybackslash}p{3cm}>{\centering\arraybackslash}p{3cm}>{\centering\arraybackslash}p{3cm}>{\centering\arraybackslash}p{3cm}}
\hline
\textbf{级别} & \textbf{特别重大（I级）} & \textbf{重大（II级）} & \textbf{较大（III级）} & \textbf{一般（IV级）} \\
\hline
\textbf{颜色} & 红色 & 橙色 & 黄色 & 蓝色 \\
\hline
\end{tabularx}
\end{table}
\end{keybox}

\subsection{应对目标与原则}

\begin{keybox}
\textbf{应对目标：}有效预防、及时控制和消除突发公共卫生事件及其危害，最大程度减少突发公共卫生事件造成的危害，保障师生员工身心健康与生命安全。

\textbf{应对原则（6条）：}

\begin{enumerate}[leftmargin=*]
    \item \imp{以人为本，减少危害}
    \item \imp{居安思危，预防为主}
    \item \imp{统一领导，分级负责}
    \item \imp{依法规范，加强管理}
    \item \imp{快速反应，协同应对}
    \item \imp{依靠科技，提高素质}
\end{enumerate}
\end{keybox}

\subsection{各方应对措施}

\begin{keybox}
\textbf{1. 政府：}
\begin{enumerate}[leftmargin=*]
    \item 建立应急组织和指挥机构
    \item 建立专家咨询委员会
    \item 建立应急处理专业技术机构
\end{enumerate}

\textbf{2. 教育行政部门：}
\begin{enumerate}[leftmargin=*]
    \item 建立组织机构
    \item 制订应急预案
    \item 应急反应
    \item 信息发布
\end{enumerate}

\textbf{3. 学校：}
\begin{enumerate}[leftmargin=*]
    \item 加强领导，实行单位领导负责制
    \item 健全学校各项健康管理制度，加强学生健康管理
    \item 认真落实突发公共卫生事件报告管理制度，确保报告信息畅通
    \item 配合卫生部门采取有效措施，及时控制事件发展蔓延
    \item 加强健康教育，提高学生应对能力
\end{enumerate}
\end{keybox}

\section{学校传染病事件}

\begin{defbox}
\textbf{学校传染病事件（infectious disease events in school）：}学校某种传染病在短时间内发生、波及范围广泛，出现大量病人或死亡病例，发病率远超常年发病率水平的情况。
\end{defbox}

\subsection{流行病学特征}

\begin{keybox}
学校突发公共卫生事件\textbf{最主要类型}。

\begin{enumerate}[leftmargin=*]
    \item \textbf{地区：}社会经济相对落后的西南省份、乡村中小学和县城小学多发。
    \item \textbf{时间：}双峰型，3-6月、10-12月。
    \item \textbf{病种：}呼吸道（流感、水痘、风疹、麻疹、流行病腮腺炎）和消化道传染病为主。
\end{enumerate}
\end{keybox}

\subsection{预防与应急管理}

\begin{keybox}
\textbf{预防：}

\begin{enumerate}[leftmargin=*]
    \item 改善学校卫生环境和教学卫生条件
    \item 提供合乎卫生要求的生活条件
    \item 开展健康教育
    \item 贯彻落实学校各项卫生管理制度
    \item 加强预警预测
\end{enumerate}

\textbf{应急管理（7步）：}

\begin{enumerate}[leftmargin=*]
    \item \imp{疫情报告}（24h内）
    \item \imp{疫情调查}
    \item \imp{控制疫情}
    \item \imp{信息发布}
    \item \imp{健康教育}
    \item \imp{应急反应终止}
    \item \imp{评估总结}
\end{enumerate}
\end{keybox}

\section{学校食物中毒}

\subsection{基本概念}

\begin{defbox}
\textbf{食源性疾病（food origin disease）：}食品中致病因素进入人体引起的感染性、中毒性疾病。3个基本要素：（1）致病因子：病毒、细菌、寄生虫、毒素、重金属、有毒有害化学物质；（2）传播媒介：食物；（3）临床表现/症状体征：轻微胃肠炎→致命神经毒作用、肝肾综合征等各不相同。

\textbf{食物中毒（food poisoning）：}食用被有毒有害物质污染的食品或食用含有毒有害物质的食品后出现的急性、亚急性疾病。特点：（1）发病与特定食物有关；（2）潜伏期短，发病曲线单峰型；（3）临床表现相似：最常见为消化道症状，如腹痛、腹泻、恶心、呕吐；（4）一般没有人与人间的直接传染。

\textbf{学校食物中毒事故（food poisoning incidents in school）：}由学校主办或管理的校内供餐单位及学校负责组织提供的集体用餐导致的学校师生食物中毒事故，分为重大、较大和一般学校食物中毒事故。
\end{defbox}

\subsection{类型}

\begin{keybox}
\textbf{食物中毒四大类型：}

\begin{enumerate}[leftmargin=*]
    \item \imp{细菌性食物中毒：}因摄入被致病菌或其毒素污染的食品后发生的急性或亚急性疾病。是最常见的食物中毒，也是学校食物中毒人数最多的类型。
    \item \imp{真菌性食物中毒：}因摄入被真菌有毒代谢产物污染的食品后所致的中毒。
    \item \imp{有毒动植物中毒：}动植物本身含某种天然有毒成分，或由于储存不当产生某种有毒物质，被人食用后所致的中毒。
    \item \imp{化学性食物中毒：}食物被化学物质污染或误食化学物质所致的中毒。
\end{enumerate}
\end{keybox}

\subsection{流行病学特征}

\begin{keybox}
学校突发公共卫生事件主要类型，仅次于传染病。

\begin{enumerate}[leftmargin=*]
    \item \textbf{地区：}南方>北方，华东、华南最高发。
    \item \textbf{时间：}秋季最高，夏季次之，冬季最少。
    \item \textbf{人群：}食物中毒起数、中毒人数：中学>小学；死亡病例集中在小学，尤其是农村小学。
    \item \textbf{病因：}食物污染和变质是最常见原因，微生物是最主要致病因子，化学物是引起死亡的主要因素。
\end{enumerate}
\end{keybox}

\subsection{预防与应急管理}

\begin{keybox}
\textbf{预防：}

\textbf{1. 学校：}食堂建筑、设备与环境；食品采购与储存；食品加工；食堂从业人员；监督与管理。

\textbf{2. 教育部门：}加强行政管理；组织人员培训；定期检查督促。

\textbf{3. 卫生和计划生育管理部门：}加强卫生监督；加大卫生许可工作管理和督查力度。

\textbf{应急管理：}

\begin{enumerate}[leftmargin=*]
    \item \imp{报告}（2h内）
    \item \imp{封存、销毁}
    \item \imp{迅速组织救治}
    \item \imp{配合调查和善后}
    \item \imp{尽早恢复正常教学秩序}
    \item \imp{正常发布相关信息}
\end{enumerate}
\end{keybox}

\begin{exambox}
\textbf{△ 报告时间对比：}传染病事件24h内报告；食物中毒2h内报告。
\end{exambox}

\section{学校群体心因性反应事件}

\subsection{概念与特征}

\begin{defbox}
\textbf{群体心因性反应（mass psychogenic reaction）：}一种群体精神性反应，是在一定社会文化背景条件下，在两人或两人以上群体中发生、有躯体性疾病症候群、但没有可检测出的器质性变化的病症。

\textbf{触发因素：}公共卫生事件为主。

\textbf{个体易患因素：}
\begin{enumerate}[leftmargin=*]
    \item 农村多发
    \item 小学和初中多发
    \item 女生>男生
    \item 认知能力较差，易接受心理暗示
\end{enumerate}
\end{defbox}

\subsection{早期识别}

\begin{keybox}
群体心因性反应事件早期原因不明，极易在学生和家长中引起恐慌，因此需重视群体心因性反应事件的\textbf{早期识别}。
\end{keybox}

\subsection{应急管理}

\begin{keybox}
\textbf{1. 应急预案（流行病学调查步骤）：}

\begin{enumerate}[leftmargin=*]
    \item \imp{描述流行病学}，形成病因假设
    \item \imp{时间分布}，确定发病高峰，同源性一次暴发/非同源性数次暴发？
    \item \imp{人群分布}
    \item \imp{地区分布}
    \item \imp{分析流行病学}，检验病因假设
    \item \imp{实验流行病学}，最终确定病因
\end{enumerate}

\textbf{2. 现场应急处置：}

\begin{enumerate}[leftmargin=*]
    \item \imp{隔离患者}
    \item \imp{消除紧张性情绪环境}
    \item \imp{对症治疗}
\end{enumerate}

\textbf{3. 心理治疗：}

\begin{enumerate}[leftmargin=*]
    \item \imp{移情解释法}
    \item \imp{暗示疗法}
    \item \imp{着重治疗关键患者（首发患者）}
    \item \imp{争取家长配合}
\end{enumerate}
\end{keybox}

% ----- 复习题 -----
\reviewcontent{
\section{复习题}

\begin{enumerate}[leftmargin=*, itemsep=8pt]
    \item \textbf{【名词解释】}学校突发公共卫生事件；学校传染病事件；食源性疾病；食物中毒；学校食物中毒事故；群体心因性反应
    \item \textbf{【简答题】}简述学校突发公共卫生事件的7种类型和4级分级。
    \item \textbf{【简答题】}简述学校突发公共卫生事件的应对目标和原则。
    \item \textbf{【简答题】}简述学校在突发公共卫生事件应对中的主要措施。
    \item \textbf{【简答题】}简述学校传染病事件的流行病学特征和应急管理步骤。
    \item \textbf{【简答题】}简述食物中毒四大类型及学校食物中毒的流行病学特征。
    \item \textbf{【简答题】}简述学校食物中毒的预防措施和应急管理要点。
    \item \textbf{【简答题】}简述群体心因性反应事件的特征、早期识别重要性和应急管理措施。
\end{enumerate}

\subsection*{参考答案要点}

\begin{enumerate}[leftmargin=*, itemsep=5pt]
    \item \textbf{学校突发公共卫生事件}：在学校内突然发生，造成或可能造成师生员工健康严重损害的重大传染病疫情、群体性不明原因疾病、重大食物或职业中毒及其他严重影响师生员工身心健康的公共卫生事件。\textbf{食源性疾病}：食品中致病因素进入人体引起的感染性、中毒性疾病。\textbf{食物中毒}：食用被有毒有害物质污染的食品后出现的急性、亚急性疾病，潜伏期短、发病曲线单峰型、临床表现相似、一般无直接人传人。\textbf{群体心因性反应}：在两人或两人以上群体中发生、有躯体性疾病症候群、但没有可检测出的器质性变化的病症。

    \item 7种类型：(1)重大传染病疫情；(2)预防接种和预防服药群体性不良反应；(3)群体性不明原因疾病；(4)食物中毒；(5)其他中毒；(6)环境因素事件；(7)意外辐射照射事件。4级：特别重大（I级）红色→重大（II级）橙色→较大（III级）黄色→一般（IV级）蓝色。

    \item 目标：有效预防、及时控制和消除突发公共卫生事件及其危害，最大程度减少危害，保障师生员工身心健康与生命安全。6条原则：以人为本减少危害；居安思危预防为主；统一领导分级负责；依法规范加强管理；快速反应协同应对；依靠科技提高素质。

    \item (1)加强领导，实行单位领导负责制；(2)健全学校各项健康管理制度，加强学生健康管理；(3)认真落实突发公共卫生事件报告管理制度，确保报告信息畅通；(4)配合卫生部门采取有效措施，及时控制事件发展蔓延；(5)加强健康教育，提高学生应对能力。

    \item 流行病学特征：学校突发公共卫生事件最主要类型；西南省份、乡村中小学和县城小学多发；双峰型（3-6月、10-12月）；以呼吸道和消化道传染病为主。应急管理7步：疫情报告(24h内)→疫情调查→控制疫情→信息发布→健康教育→应急反应终止→评估总结。

    \item 四大类型：(1)细菌性（最常见，人数最多）；(2)真菌性；(3)有毒动植物中毒；(4)化学性。流行病学特征：南方>北方，华东华南最高发；秋季最高，夏季次之；中毒起数中学>小学，死亡病例集中在小学特别是农村小学；食物污染和变质最常见原因，微生物最主要致病因子，化学物是引起死亡的主要因素。

    \item 预防：(1)学校：食堂建筑设备环境、食品采购储存、食品加工、从业人员、监督管理；(2)教育部门：行政管理、人员培训、定期检查；(3)卫生部门：卫生监督、许可管理和督查。应急：(1)报告(2h内)；(2)封存销毁；(3)迅速组织救治；(4)配合调查善后；(5)尽早恢复正常教学秩序；(6)正常发布相关信息。

    \item 特征：群体精神性反应，有躯体症候群但无器质性变化，触发因素以公共卫生事件为主，农村>城市，小学初中多发，女生>男生，认知能力较差者易感。早期识别：早期原因不明极易引起恐慌，需高度重视。应急管理：(1)应急预案：描述流行病学→时间/人群/地区分布→分析流行病学→实验流行病学；(2)现场处置：隔离患者、消除紧张情绪环境、对症治疗；(3)心理治疗：移情解释法、暗示疗法、着重治疗首发患者、争取家长配合。
\end{enumerate}
}

\newpage
